\documentclass[11pt,a4paper]{article}
\usepackage[margin=28mm]{geometry}
\usepackage{amsmath,amssymb,amsthm,mathtools}
\usepackage[T1]{fontenc}
\usepackage{lmodern,microtype}
\usepackage{xcolor,enumitem}
\usepackage[colorlinks=true,linkcolor=blue!45!black,citecolor=blue!45!black,urlcolor=blue!45!black]{hyperref}
\usepackage{fancyhdr}
\setlength{\headheight}{14pt}
\pagestyle{fancy}
\fancyhf{}
\fancyhead[L]{\small An explicit three-adic Gersten obstruction}
\fancyhead[R]{\small Research draft}
\fancyfoot[C]{\thepage}
\setlength{\parskip}{3pt}
\setlength{\parindent}{14pt}
\setlist{itemsep=3pt,topsep=5pt}
\makeatletter
\renewcommand*\l@section[2]{\vskip0.25em\@dottedtocline{1}{0em}{1.8em}{\bfseries #1}{#2}}
\makeatother
\newtheorem{theorem}{Theorem}[section]
\newtheorem{proposition}[theorem]{Proposition}
\newtheorem{lemma}[theorem]{Lemma}
\newtheorem{corollary}[theorem]{Corollary}
\theoremstyle{remark}
\newtheorem{remark}[theorem]{Remark}
\newcommand{\Q}{\mathbf Q}
\newcommand{\Z}{\mathbf Z}
\newcommand{\R}{\mathbf R}
\newcommand{\C}{\mathbf C}
\newcommand{\F}{\mathbf F}
\newcommand{\A}{\mathbf A}
\newcommand{\PP}{\mathbf P}
\newcommand{\OO}{\mathcal O}
\newcommand{\m}{\mathfrak m}
\newcommand{\n}{\mathfrak n}
\newcommand{\Spec}{\operatorname{Spec}}
\newcommand{\Pic}{\operatorname{Pic}}
\newcommand{\Cl}{\operatorname{Cl}}
\newcommand{\Frac}{\operatorname{Frac}}
\newcommand{\divisor}{\operatorname{div}}
\newcommand{\colim}{\operatorname*{colim}}
\newcommand{\HD}{H_{\mathcal D,\mathrm{an}}}
\newcommand{\reg}{\operatorname{reg}}
\newcommand{\Qrat}{_{\Q}}
\newcommand{\Gal}{\operatorname{Gal}}
\title{\textbf{An explicit three-adic obstruction\newline to rational Gersten injectivity}}
\author{Research draft prepared with ChatGPT}
\date{7 October 2026}
\begin{document}
\maketitle
\thispagestyle{empty}
\begin{center}
\fcolorbox{black!30}{black!3}{\parbox{0.9\textwidth}{\small
\textbf{Review status.} This is a proposed new counterexample and proof,
prepared for independent mathematical review. The argument has been
internally checked, including adversarial checks, but has not been
externally refereed or formally verified. The cited papers supply the
standard constructions; they do not assert the new counterexample.}}
\end{center}
\begin{abstract}
We propose an explicit non-torsion class in the generic kernel of
Quillen $K_3$ for the three-dimensional regular local ring
$(\Z_{(3)}[x,y,z]/(3+xy+z^2))_{(x,y,z)}$ and for its completion.
The class is twice a product of three units. Dennis--Stein identities
make its image zero after inverting $x$, integrally. To detect infinite
order, we map the completion to a smooth formal quadric over
$\Z_3[\sqrt{-3}]$ and calculate a relative continuous cyclic character
using the Beilinson fiber square of Antieau--Mathew--Morrow--Nikolaus.
A two-chart residue calculation reduces nonvanishing to an elementary
convergent three-adic series, which is $4$ modulo $3$ after normalization.
The final injection is between actual relative and absolute $K$-groups,
using $K_4$ of the special fiber equal to zero. This short proof uses
neither a number-field regulator nor a comparison of syntomic
normalizations.
\end{abstract}

\section{The ring and the explicit class}
All $K$-groups below are actual Quillen algebraic $K$-groups unless a
completed spectrum is explicitly indicated. Put
\[
 A=\left(\Z_{(3)}[x,y,z]/(3+xy+z^2)\right)_{(x,y,z)},
 \qquad
 \widehat A=\Z_3[[x,y,z]]/(3+xy+z^2).
\]
Both rings are regular local of dimension three: the ambient regular
local ring has dimension four, and the defining relation has nonzero
linear term $3$ in its maximal ideal modulo its square. Equivalently,
the quotient has dimension three and maximal ideal generated by
$x,y,z$. In particular each is a domain. They are ramified in the
sense that $3=-xy-z^2$ belongs to the square of the maximal ideal.

\begin{theorem}[Proposed explicit counterexample]\label{thm:shortmain}
For $R=A$ or $\widehat A$, the element
\[
 \boxed{\quad c_R=2\left\{\frac{1+z}{2},\,1+x,\,
                 1+\frac{(1+x)y}{4}\right\}\in K_3(R)\quad}
\]
has infinite order, while its image in $K_3(R[1/x])$ is zero.
Consequently the generic map
$K_3(R)_{\Q}\to K_3(\Frac R)_{\Q}$ has nonzero kernel.
\end{theorem}
Here braces mean the product of the three unit classes in Quillen
$K_1$. Each entry is a unit in both local rings. The proof will establish
nonvanishing after mapping $\widehat A$ to an auxiliary ring, and
therefore proves it for $A$ as well.

\begin{lemma}\label{lem:shortgenericzero}
The displayed class vanishes in $K_3(R[1/x])$ integrally.
\end{lemma}
\begin{proof}
Put $a=(1+z)/2$, $t=y/4$, $q=a(1-a)=1+xt$,
$\lambda=1+x$, and $\mu=1+\lambda t$. All these expressions,
including $1-a$, are units when required below. Use the
Dennis--Stein convention in which $\langle r,s\rangle$ is defined
when $1-rs$ is a unit. Its elementary relations give
\[
 \delta=\langle-x,t\rangle
       =\langle-x,\mu\rangle=\{\lambda q,\mu\}.
\]
Indeed $\langle-x,1\rangle=0$, the additivity relation replaces
$t$ by $t+1+xt=\mu$, and $1+x\mu=\lambda q$.
See \cite[III, Definition~5.11 and (D1)--(D3)]{Weibel} for these
relations and conventions. Steinberg and graded commutativity give
\[
 2\{a,q\}=2\{a,a\}+2\{a,1-a\}=0.
\]
Hence $c_R=2[a]\delta$. After inverting $x$, the same
Dennis--Stein symbol equals $\{-x,q\}$, and
\[
 c_R=2\{a,-x,q\}=-2\{a,q,-x\}=0.
\]
No rational coefficient is used in this vanishing argument.
\end{proof}

\section{A relative character and its unit formula}
Set $W=\Z_3[u]/(u^2+3)$, $E=W[1/3]$, and $\pi=u$.
The ring $W$ is a complete DVR. Write
$K(S;\Q_3)=K(S)^{\wedge}_3[1/3]$ for the rationalization of the
three-completed \emph{spectrum}; this does not denote
$K_*(S)\otimes\Q_3$ on homotopy groups.
Continuous cyclic invariants and continuous differential forms are
completed before inverting $3$. We use the ordinary derived category
and the ordinary quotient by the image of a differential.
In particular no image is replaced by its topological closure.
Write $\Omega_B^j=\widehat\Omega^j_{B/W}[1/3]$.

\begin{lemma}[Relative character of a product of units]
\label{lem:relativelogproduct}
Suppose that $\mathfrak X=\operatorname{Spf}B$ is smooth, affine,
topologically of finite presentation, and of relative dimension at
most two over $\operatorname{Spf}W$. There is a natural relative map
\[
 \chi_B:K(B,\pi B)\longrightarrow
                  \Sigma HC^{\mathrm{cts}}(B/W;\mathbf Q_3).
\]
For $v\in1+\pi B$ and $b_1,b_2\in B^\times$, the top HKR component
of its degree-three value is
\begin{equation}\label{eq:tcreviewlogproduct}
 \chi_B([v][b_1][b_2])_{\mathrm{top}}
   =\log(v)\,d\log b_1\wedge d\log b_2
       \quad\text{in }\Omega_B^2/d\Omega_B^1,
\end{equation}
up to one universal sign determined by the norm and suspension
conventions. Here $[v]$ has its relative lift furnished by its reduction
to the identity unit.
\end{lemma}

\begin{proof}
The fiber square of \cite[Proposition~4.10]{AMMN} has rows
\[
\begin{matrix}
 K^{\mathrm{cts}}(\mathfrak X;\mathbf Q_3)
       &\longrightarrow&K(B/\pi;\mathbf Q_3)\\
 \downarrow&&\downarrow\\
 TC^{\mathrm{cts}}(\mathfrak X;\mathbf Q_3)
       &\longrightarrow&TC(B/\pi;\mathbf Q_3)\\
 \downarrow&&\downarrow\\
 HC^{-,\mathrm{cts}}(B/W;\mathbf Q_3)
       &\longrightarrow&HP^{\mathrm{cts}}(B/W;\mathbf Q_3).
\end{matrix}
\]
The last horizontal fiber is $\Sigma HC^{\mathrm{cts}}(B/W;\mathbf Q_3)$,
not negative cyclic homology. Mapping the actual relative spectrum to
these fibers gives $\chi_B$. Consequently relative $K_3$ has target
$HC_2$, and its top component is $\Omega_B^2/d\Omega_B^1$.

We record why the fiber comparison respects multiplication by absolute
$K$-classes. Set $H=THH(B)$, a commutative algebra in cyclotomic spectra.
If $TC_{\mathrm{cyc}}$ denotes the functor on that category, then
$TC_{\mathrm{cyc}}(H)=TC(B)$; no second $THH$ is intended.
All objects $H\otimes D$ and maps $\mathrm{id}_H\otimes f$ in
\cite[Proposition~2.8, Corollary~2.10, Theorem~2.12]{AMMN} are
$H$-modules and module maps. The fixed-point and Tate functors, their
canonical maps, and the norm triangle preserve the induced $TC(B)$-actions
\cite[Corollary~I.4.3]{NikolausScholze}. Thus the inverses of the
equivalences used there are module inverses. The comparisons from
derived to ordinary reduction and then to $B/\pi$ are induced by
$B$-algebra maps. The projection $HH(B)\to HH(B/W)$ in the proof of
Proposition~4.10 also preserves the action. These inverses and fiber
equivalences are used after $3$-completion and rationalization.

The multiplicative cyclotomic trace
\cite[Theorem~7.4]{BGTTrace} now gives $K(B)$-linearity, with target
action through the canonical map $K(B)\to HC^{-,\mathrm{cts}}(B/W;\mathbf Q_3)$.
At the cyclic-complex level, the compatibility of the norm boundary
and its suspension with the $HC^-$-module action is the Connes-sequence
compatibility of \cite[Proposition~2.3]{BokstedtOttosen}; the spectral norm is the same Connes operator by \cite[Theorem~2.1, Lemma~2.2, and Theorem~2.3]{HoyoisCircle}. This includes
the usual graded sign of a connecting map. In particular, no assumption
that an unspecified additive fiber equivalence is multiplicative is
being made.

To normalize a relative unit, first take $B=W\langle T\rangle$ and
$v=1+\pi T$. Its relative character is a function
$f\in HC_0^{\mathrm{cts}}(B/W;\mathbf Q_3)=E\langle T\rangle$.
Choose the norm convention in which its image in $HC^-_1$ is $df$.
The left vertical arrow in \cite[Theorem~2.12]{AMMN} is canonical;
its projection to Hochschild homology takes a unit to the Dennis trace
$v^{-1}\otimes v$. Hence $df=d\log v$. Evaluation at $T=0$ gives
$f(0)=0$. Continuous differentiation on $E\langle T\rangle$ has
only constants in its kernel, and therefore
\[
 f=\log(1+\pi T)=\sum_{n\ge1}(-1)^{n+1}\pi^nT^n/n.
\]
The series converges beyond the closed unit disc. Substitution
$T=(v-1)/\pi$ proves the assertion for every relative unit under
consideration.

Finally use the normalized cyclic complex
$C^+=E[t,t^{-1}]/tE[t]\otimes C_*$, with $|t|=-2$.
The $HC^-\times HC$ action is shuffle plus $t$ times cyclic shuffle
\cite[Definition~2.1 and Remark~2.2]{BokstedtOttosen}.
Starting with the column-zero representative $\log v$, all positive
columns and cyclic-shuffle corrections vanish in $C^+$. After each
of the two unit multiplications only the ordinary shuffle remains.
Under
\[
 b_0\otimes\cdots\otimes b_n\longmapsto
       b_0\,db_1\wedge\cdots\wedge db_n/n!,
\]
the two shuffles in degree two cancel the factor $2!$, giving
\eqref{eq:tcreviewlogproduct}, with the suspension signs just specified.
This calculation involves only degrees through two and commutes with
the continuous, relative-to-$W$ projection. It does not assume a
multiplicative splitting of an infinite HKR filtration.
\end{proof}


\paragraph{The finite continuous HKR calculation.}
For the three-adic quadric, put
\[
 B_0=W[X,Y,Z]/(XY+Z^2-1),\qquad B=\widehat{B_0}_3,
 \qquad \Omega_B^i=\widehat\Omega^i_{B/W}[1/3].
\]
The algebra $B_0$ is smooth over $W$ of relative dimension two.
The normalized HKR map on its Hochschild complex is
\[
 b_0\otimes\cdots\otimes b_n\longmapsto
 \frac{1}{n!}b_0\,db_1\wedge\cdots\wedge db_n
 \quad(0\le n\le2),
\]
and is zero in higher degrees.  It is defined over $W$, since $2$ is
a unit, and intertwines the Hochschild differential with zero and the
Connes operator with $d$.  Smoothness makes it a quasi-isomorphism of
mixed complexes: the Hochschild homology is $\Omega^n_{B_0/W}$ for
$0\le n\le2$ and vanishes above degree two.
The comparison with the continuous invariant is supplied by
\cite[Proposition~4.2]{AMMN}.  In the cyclic, negative cyclic, and
periodic total complexes of this finite mixed complex, each total
degree has only finitely many differential-form terms.  Their
reduction towers are surjective, so derived three-adic completion is
computed termwise.  After inverting $3$, this gives in particular
\[
 HC^{\mathrm{cts}}_0(B/W;\mathbf Q_3)=\Omega_B^0,
 \qquad
 HC^{\mathrm{cts}}_2(B/W;\mathbf Q_3)
   =\frac{\Omega_B^2}{d\Omega_B^1}
     \oplus\ker(d:\Omega_B^0\to\Omega_B^1).
\]
The norm map from the degree-zero cyclic term to the degree-one
negative cyclic term is $f\mapsto df$, with the chosen global
suspension sign.  In particular, the displayed quotient is the
ordinary quotient by the image of $d$; its image is not replaced by
its closure.  The same calculation in relative dimension one applies
to the universal unit in $W\langle T\rangle$.



\section{A residue obstruction on the quadric}
\begin{lemma}[A residue obstruction on the good quadric]
\label{lem:tcreviewresidue}
Let $a_0,a_1\in W$ have unit difference, and put
\[
 B=W\langle X,Y,Z\rangle/
           \bigl(XY+(Z-a_0)(Z-a_1)\bigr).
\]
For $F\in E\langle Z\rangle$, the form
$\omega_F=(dX/X)\wedge dF$ is globally regular. If it is exact in
$\Omega_B^2/d\Omega_B^1$, then $F(a_0)=F(a_1)$.
\end{lemma}
\begin{proof}
The two root-deletion affine opens have completed coordinates
\[
 Z=a_i+Xt_i,\qquad
 Y=-t_i(a_i-a_{1-i}+Xt_i)\quad(i=0,1).
\]
They are the completions of $D(X)\cup D(Z-a_{1-i})$, cover the formal
scheme, and have intersection $W\langle X,X^{-1},Z\rangle$.
The forms
\[
 q_i=-(F(Z)-F(a_i))\,dX/X
\]
are regular on the respective bidiscs, satisfy $dq_i=\omega_F$, and
have difference $(F(a_1)-F(a_0))dX/X$.

A closed bidisc form $A(X,t)dX+B(X,t)dt$ has zero Laurent residue
after $t=c/X$, for every $|c|\le1$. Indeed its residue is the
convergent series
\[
 \sum_{m\ge0}(a_{m,m+1}-b_{m+1,m})c^{m+1}=0,
\]
because closedness equates $(m+1)a_{m,m+1}$ and
$(m+1)b_{m+1,m}$. Division by $m+1$ is used only in this individual
coefficient equality, not to construct a bounded primitive.
If $\omega_F=dq$ globally, apply this to $q_i-q$ after restricting
the overlap to $Z=0$, where $c=-a_i$. The difference then has residue
zero, which is the required endpoint equality. This argument also includes $c=0$.
\end{proof}


\section{The explicit endpoint calculation}
\begin{proof}[Proof of Theorem~\ref{thm:shortmain}]
With $W,E,u$ as fixed above, put
\[
 B=W\langle X,Y,Z\rangle/(XY+Z^2-1),\qquad C=B/u.
\]
This is a smooth affine formal $W$-scheme of relative dimension two,
topologically of finite presentation. Substitution
$(x,y,z)=(uX,uY,uZ)$ defines an actual ring map $\widehat A\to B$, since all
formal series converge. In $B$ write
\[
 a=\frac{1+uZ}{2},\quad \lambda=1+uX,\quad
 \mu=1+\frac{(1+uX)uY}{4},\quad
 q=a(1-a)=\frac{1+3Z^2}{4}.
\]
These are units, $a^2\in1+uB$, and the image of $c$ has relative lift
$[a^2][\lambda][\mu]$.

The ruling presents $\operatorname{Spec}C$ as an affine-line torsor
over $\mathbf P^1_{\mathbf F_3}$. Its two base charts are
\[
 Z=1+Xt,\quad Y=-t(2+Xt),\qquad
 Z=-1-sY,\quad X=-s(2+sY),\quad s=t^{-1}.
\]
Thus $K(C)\simeq K(\mathbf F_3)\oplus K(\mathbf F_3)$ by homotopy
invariance, descent, and the projective bundle formula. Quillen's
finite-field computation gives $K_4(C)=0$ for the actual groups.
The exact sequence of the actual relative spectrum therefore gives
\[
 0=K_4(C)\longrightarrow K_3(B,uB)\longrightarrow K_3(B),
\]
so the last arrow is injective. It suffices to find a nonzero value
of the relative character in its rational top-form component:
this makes the specified actual relative lift, and then its actual
absolute image, have infinite order.

Set
\[
 L=-\log2+\log(1+uZ),\qquad
 \omega=\frac{dX}{X}\wedge d\log q.
\]
Lemma~\ref{lem:relativelogproduct} gives, up to its fixed sign, the
global form $\tau=2L\,d\log\lambda\wedge d\log\mu$.
Direct differentiation gives
\[
 d\log\lambda\wedge d\log\mu
   =\frac{dX}{X}\wedge d\log q-d\log q\wedge d\log\mu.
\]
Choose $F\in E\langle Z\rangle$ with
\[
 dF=L\,d\log q=\frac{6ZL(Z)}{1+3Z^2}\,dZ.
\]
The right side converges on every closed disc of radius
$1<R<\sqrt3$; termwise integration still has radius greater than one.
There is an equality of global forms
\begin{equation}\label{eq:directtcglobalform}
 \tau=2\frac{dX}{X}\wedge dF-2d(F\,d\log\mu).
\end{equation}
Lemma~\ref{lem:tcreviewresidue}, with roots $1,-1$, shows that its
class is nonzero if $F(1)-F(-1)\ne0$.

The odd part of $L$ is
$u\sum_{k\ge0}(-3)^kZ^{2k+1}/(2k+1)$. Expanding the denominator
and integrating its even product with $6Z$ gives
\[
 \frac{F(1)-F(-1)}u
     =12\sum_{j\ge0}\frac{(-3)^jS_j}{2j+3},
 \qquad S_j=\sum_{k=0}^j\frac1{2k+1}.
\]
The first summand is $4$. For $j\ge1$, the valuation of the $j$th
summand is at least
\[
 A_j=1+j-\lfloor\log_3(2j+1)\rfloor-v_3(2j+3).
\]
First, $A_j\ge1$: for $j=1,2,3$ the lower bounds are respectively
$1,2,1$; for $j\ge4$ use
$\lfloor\log_3(2j+3)\rfloor\le\lfloor j/2\rfloor$.
The latter follows from the cases $j=4,5$ and
$2j+7<3(2j+3)$, advancing $j$ by two. Separately,
\[
 A_j\ge1+j-2\log_3(2j+3)\longrightarrow+\infty.
\]
This proves convergence of the tail as well as its divisibility by
$3$. Hence $(F(1)-F(-1))/u\in4+3\mathbf Z_3$, which is nonzero.
The character of the actual relative lift is nonzero in a
$\mathbf Q_3$-vector space. The actual relative-to-absolute injection
just proved shows that its image $c_B\in K_3(B)$ has infinite order.
Consequently the original actual element $c\in K_3(\widehat A)$ has infinite
order as well. Completion constructs the character; it is not needed
in this last injectivity assertion.
\end{proof}
\section{The ordinary target and the scope of the claim}
\paragraph{The ordinary quotient is essential.}
The target in \cite[Constructions~4.6 and~4.11]{AMMN} uses the
ordinary derived category. For example, in the preceding quadric put
$\omega_0=(dX/X)\wedge dZ$. The globally regular primitives
$-(Z-Z^{3^j})dX/X$ satisfy
\[
 d\!\left(-(Z-Z^{3^j})\frac{dX}{X}\right)
   =\omega_0-3^jZ^{3^j-1}\omega_0\longrightarrow\omega_0.
\]
The ruling form is therefore in the closure of exact forms although
the residue lemma shows it is not exact. No assertion about
Hausdorff-reduced cohomology or the order at each finite Artin stage
is made here.


The same formulas prove the assertion for the henselization $A^h$:
it is regular local, maps to $\widehat A$, and the integral vanishing
identity already holds after inverting $x$ in $A^h$.
The three-dimensional example is enough to contradict the first
injectivity requirement of the unrestricted augmented rational
Gersten complex, if the proposed argument is correct.

\paragraph{The precise point for independent review.}
The series estimate and the residue obstruction are elementary.
The main interface requiring a specialist's review is the
$K(B)$-module compatibility of the AMMN horizontal-fiber comparison,
including the normalization of a relative unit and its identification
with the classical cyclic norm. Those steps are stated with their
constructions and primary references above. No claim of external
verification or priority is made here.

\begin{thebibliography}{9}
\small
\setlength{\itemsep}{2pt}
\bibitem{AMMN}
B.~Antieau, A.~Mathew, M.~Morrow, and T.~Nikolaus,
\emph{On the Beilinson fiber square}. Propositions~2.8, 4.10, 4.12,
Corollary~2.10, Theorem~2.12, and Constructions~4.6, 4.11.
\href{https://arxiv.org/abs/2003.12541}{arXiv:2003.12541}.

\bibitem{BGTTrace}
A.~J.~Blumberg, D.~Gepner, and G.~Tabuada,
\emph{Uniqueness of the multiplicative cyclotomic trace},
Advances in Mathematics \textbf{260} (2014), 191--232, Theorem~7.4.
\href{https://arxiv.org/abs/1103.3923}{arXiv:1103.3923}.

\bibitem{NikolausScholze}
T.~Nikolaus and P.~Scholze, \emph{On topological cyclic homology},
Acta Mathematica \textbf{221} (2018), 203--409, Corollary~I.4.3.
\href{https://arxiv.org/abs/1707.01799}{arXiv:1707.01799}.

\bibitem{BokstedtOttosen}
M.~B\"okstedt and I.~Ottosen,
\emph{A Hochschild--Kostant--Rosenberg theorem for cyclic homology},
Section~2, especially Definition~2.1, Remark~2.2, and Proposition~2.3.
\href{https://arxiv.org/abs/1601.07412}{arXiv:1601.07412}.

\bibitem{Quillen}
D.~Quillen, \emph{Higher algebraic $K$-theory I}, in \emph{Algebraic $K$-Theory I: Higher $K$-Theories}, Lecture Notes in Mathematics \textbf{341}, Springer, 1973, 85--147.
\href{https://doi.org/10.1007/BFb0067053}{doi:10.1007/BFb0067053}.

\bibitem{Weibel}
C.~A.~Weibel, \emph{The $K$-book: An Introduction to Algebraic $K$-theory}, Graduate Studies in Mathematics \textbf{145}, American Mathematical Society, 2013. Chapter~III, Definition~5.11 and relations (D1)--(D3) fix the Dennis--Stein convention used here.
\href{https://sites.math.rutgers.edu/~weibel/Kbook/Kbook.III.pdf}{Author's Chapter~III}.
\bibitem{HoyoisCircle}
M.~Hoyois, \emph{The homotopy fixed points of the circle action on
Hochschild homology}, 2018. Theorem~2.1, Lemma~2.2, and Theorem~2.3.
\href{https://arxiv.org/abs/1506.07123}{arXiv:1506.07123}.

\end{thebibliography}
\end{document}
